\chapter{Introducción y Objetivos}
\label{chap:introduccion}

La configuración en emisor común con transistor bipolar de juntura (BJT) constituye una de las topologías
amplificadoras fundamentales en la electrónica analógica discreta. Su amplia difusión se debe a su capacidad
de proveer simultáneamente una ganancia de tensión significativa y una ganancia de corriente apreciable,
brindando una ganancia de potencia global elevada con respecto a otras configuraciones elementales.

\section{Objetivos del Trabajo Práctico}

El presente trabajo de laboratorio persigue los siguientes objetivos técnicos y formativos:
\begin{enumerate}
  \item Diseñar un amplificador con transistor BJT NPN en configuración emisor común polarizado por divisor resistivo,
    estableciendo el punto de operación en reposo $Q(V_{CEQ}, I_{CQ})$ bajo el criterio de Máxima Excursión Simétrica (MES).
  \item Garantizar la estabilidad del punto de reposo frente a dispersiones del parámetro $\beta$ y variaciones
    de temperatura ambiental mediante la adecuada elección de la red de polarización y resistencia de emisor.
  \item Evaluar el comportamiento del circuito mediante simulación computarizada en LTspice, analizando el punto de
    operación en continua (.op) y la respuesta temporal en régimen transitorio (.tran) ante señales senoidales de $1\kHz$.
  \item Normalizar los componentes calculados adoptando resistores comerciales de las series estándar (E12),
    verificando que la desviación respecto a los valores analíticos no exceda la tolerancia máxima admitida del $10\%$.
  \item Realizar el análisis gráfico de las rectas de carga estática de corriente continua (CC) y dinámica de corriente
    alterna (CA), cuantificando la excursión simétrica real sin distorsión sobre la carga.
  \item Determinar mediante el modelo híbrido equivalente de pequeña señal los parámetros circuitales fundamentales:
    impedancia de entrada ($Z_i$), impedancia de salida ($Z_o$), ganancia de tensión ($A_v$) y ganancia de corriente ($A_i$).
  \item Montar físicamente el circuito en placa de circuito impreso (PCB) en laboratorio y relevar de manera experimental
    las magnitudes estáticas de polarización y los parámetros dinámicos empleando el método del resistor serie ($R_S$).
  \item Contrastar de forma cuantitativa los resultados analíticos, las simulaciones y las mediciones de laboratorio.
\end{enumerate}

\section{Especificaciones de Partida y Circuito Base}

Los parámetros iniciales del circuito fueron suministrados por la cátedra en la consigna de trabajo:
\begin{itemize}
  \item Tensión continua de alimentación: $V_{CC} = 15\V$.
  \item Resistencia de emisor para polarización: $R_E = 180\ohm$.
  \item Resistencia de colector: $R_C = 1.2\kohm$.
  \item Resistencia de carga conectada a la salida: $R_L = 1\kohm$.
  \item Frecuencia de señal de prueba de alterna: $f = 1\kHz$.
  \item Transistor NPN de silicio de uso general seleccionado: BC547B, con ganancia de corriente medida de $\beta = 284$.
  \item Caída de tensión directa base-emisor estimada para silicio: $V_{BE} = 0.70\V$.
\end{itemize}

En la figura~\ref{crkt:emisor_comun} se representa el diagrama esquemático completo de la etapa amplificadora a implementar.

\begin{figure}[!ht]
  \centering
  \begin{circuitikz}[scale=0.95, transform shape]
  % Referencia de masa
  \draw (0,0) node[ground] {}
    to[sV, l=$v_s$] (0,3)
    to[R, l=$R_s$] (1.8,3)
    to[C, l=$C_1$, -*] (3.5,3) coordinate (base_node);

  % Divisor resistivo de entrada
  \draw (3.5,6.5) node[vcc] {$V_{CC}$} coordinate (vcc_node)
    to[R, l=$R_2$, -*] (base_node)
    to[R, l=$R_1$, *-] (3.5,0) node[ground] {};

  % Transistor BJT NPN
  \draw (base_node) to[short] (4.2,3) node[npn, anchor=B] (Q1) {};
  \node[right=4pt] at (Q1.right) {$Q_1$ (BC546B)};

  % Malla de colector
  \draw (Q1.C) to[short, -*] ++(0,0.25) coordinate (coll_node)
    to[R, l_=$R_C$] ++(0,2) coordinate (aux) to [short] (vcc_node |- aux);

  % Malla de emisor con capacitor de desacoplo
  \draw (Q1.E) to[short, -*] ++(0,-0.5) coordinate (emit_node);
  \draw (emit_node) to[R, l=$R_E$] (5,0) node[ground] {};
\draw (emit_node) to[short] ++(1.5,0) coordinate(aux)
      to[C, l=$C_E$] (aux |- 0,0) node[ground] {};

  % Red de acoplo de salida y carga
  \draw (coll_node) to[C, l=$C_2$] ++(4,0) coordinate (out_node)
    to[short, -o] ++(1,0) node[above] {$v_L$};
  \draw (out_node) to[R, l=$R_L$, *-] (9,0) node[ground] {};
\end{circuitikz}

  \caption{Esquema circuital completo del amplificador en configuración emisor común con desacoplo de emisor.}
  \label{crkt:emisor_comun}
\end{figure}

El capacitor $C_1$ cumple la función de desacoplar el nivel continuo de la señal de entrada $v_s$, impidiendo que la
resistencia interna de la fuente modifique la polarización de base. De forma análoga, el capacitor de acoplo de salida
$C_2$ bloquea la tensión continua presente en el colector ($V_{CQ}$), asegurando que sobre la carga $R_L$ se desarrolle
únicamente la componente senoidal amplificada $v_L$. Finalmente, el capacitor de paso o desacoplo de emisor $C_E$ actúa
como un cortocircuito para la componente alterna a la frecuencia de trabajo ($1\kHz$), derivando la señal directamente a
masa y maximizando consecuentemente la ganancia de tensión de la etapa sin perjudicar la estabilidad en continua.
